
\begin{ledgroupsized}[r]{120mm}
\footnotesize
\pstart
\noindent\textbf{Überlieferung:}
\pend
\end{ledgroupsized}
\begin{ledgroupsized}[r]{114mm}
\footnotesize
\pstart
\parindent -6mm
\makebox[6mm][l]{\textit{L}}%
Notiz:
LH~XXXV~10, 8~Bl.~14. 
Ein Zettel (ca.~4,5~x~11,5~cm.);
alle Ränder bis auf den oberen beschnitten. 
Eine Seite auf Bl.~14~r\textsuperscript{o}; Bl.~14~v\textsuperscript{o} leer.
\pend
\end{ledgroupsized}
\vspace{5mm}
\begin{ledgroup}
\footnotesize
\pstart
\noindent
\textbf{Datierungsgründe:}
Die kurze Notiz N.\;\ref{60066} hat Leibniz auf einem kleinen Papierstreifen entworfen, der weder ein Wasserzeichen noch eine gemeinsame Schnittkante mit anderen bekannten Textträgern aufweist.
Die Datierung der Notiz wird zusätzlich durch die knappe Behandlung einer sehr allgeimeinen Frage erschwert: Wie nämlich aus kosmologischen Annahmen, wie etwa der Existenz eines gemeinsamen Schwerpunkts aller Körper, bewiesen werden könne, dass ein jeder Körper einen eigenen Schwerpunkt hat.
\pend
\pstart
Aus den genannten Gründen erscheint eine präzise Datierung nicht ohne Weiteres möglich. Dennoch lässt der Schreibduktus, der der (frühen) Hannoveraner Zeit als der Pariser Zeit näherkommt, die Vermutung zu, Leibniz habe die Notiz ab 1677 verfasst.
\pend 
\end{ledgroup}
\selectlanguage{latin}
\vspace{8mm}
\pstart
\centering
\href{\Scan{LH_35_10_08_014r}}%
{\lbrack14~r\textsuperscript{o}\rbrack}  
\hfill
Demonstratio centri\indexS{centrum gravitatis}\indexS{demonstratio centri gravitatis}
\edlabel{35_10_08_014_1a}%
\edtext{}{\xxref{35_10_08_014_1a}{35_10_08_014_1b}\lemma{gravitatis}\Bfootnote{%
\textit{(1)}~Fingamus %
\textit{(2)}~Si ponatur~\textit{L}}}%
gravitatis
\hfill
\phantom{\lbrack14~r\textsuperscript{o}\rbrack}
\pend 
\nobreak
\vspace{0.5\baselineskip}
\pstart
\noindent
Si ponatur\edlabel{35_10_08_014_1b}
omnia corpora tendere\indexS{corpora omnia versus centrum gravium tendentia} versus aliquod
\edtext{punctum seu centrum\indexS{centrum gravium} gravium, necesse est in ipsis esse}{%
\lemma{punctum}\Bfootnote{%
\bars{seu centrum gravium \textit{erg.}},
\textit{(1)}~est in ipsis ce
\textit{(2)}~necesse est in ipsis esse~\textit{L}}}
certum aliquod centrum\indexS{centrum gravitatis}
\edtext{gravitatis, quo cum centro}{\lemma{gravitatis,}\Bfootnote{%
\textit{(1)}~alioqui 
\textit{(a)}~punctum 
\textit{(b)}~ipsorum %
\textit{(2)}~quo cum %
\textit{(a)}~cent %
\textit{(b)}~punct %
\textit{(c)}~centro~\textit{L}}}
gravium\indexS{centrum gravium} coincidant.
\pend
\pstart
Haec demonstratio\indexS{demonstratio centri gravitatis} 
oritur ex nisu ad unum punctum.\indexS{nisus ad unum punctum}
Aliter demonstrari potest ex nisu parallelo\indexS{nisus parallelus} quanquam 
et illa accommodari possit ad hanc sumendo punctum illud valde remotum.
\pend
